%% ============================================================
%%  CHAPTER 14 — CONCLUSION AND FUTURE WORK
%% ============================================================
\chapter{Conclusion and Future Work}
\label{ch:conclusion}

This final chapter closes the report. Section~\ref{sec:conc_summary}
summarises what was built and achieved, against the objectives and
contributions set out in Chapter~1. Section~\ref{sec:conc_future}
consolidates every concrete, unresolved future-work item raised
across the report into a single organised list.
Section~\ref{sec:conc_remarks} offers closing remarks on the
project's scope and its place as a foundation for further work.

%% ----------------------------------------------------------
\section{Summary of Achievement}
\label{sec:conc_summary}
%% ----------------------------------------------------------

This project set out to build an artifact-based \acrshort{bci}
system that translates voluntary \acrshort{eog} and \acrshort{emg}
gestures --- blinks and jaw movements --- into discrete commands for
a 3D-printed prosthetic hand, evaluated under a protocol strict
enough to be trusted as a genuine estimate of deployment
performance. Measured against the seven contributions claimed in
Chapter~1, and checked against concrete evidence in Chapter~13, the
central results are:

\begin{itemize}[leftmargin=*]
    \item A \textbf{three-way hybrid classification architecture} ---
    hand-crafted features for blink/muscle routing, \acrshort{dtw}
    shape-matching for blink sub-classification, Riemannian
    covariance geometry for muscle sub-classification --- reaching
    \textbf{91.3\% accuracy} (unweighted \acrshort{loso} mean) or
    \textbf{91.8\%} (pooled), against an 85.1\% flat-classifier
    baseline (Chapter~9).
    \item A \textbf{rigorous, leakage-audited evaluation
    methodology}, including a four-test diagnostic audit and a
    label-permutation shuffle test, that both diagnosed an 8--10-point
    inflation in the naive evaluation protocol and established
    Leave-One-Session-Out as the honest reporting standard used
    throughout this report (Chapter~9).
    \item A \textbf{purpose-built, three-session dataset} of
    voluntary artifact recordings (Chapter~6), collected,
    diagnosed, and iteratively refined --- including discovering
    and fixing a silent data-loading bug and a critical single-blink
    contamination problem that directly shaped the final
    classification architecture.
    \item A \textbf{complete, physically fabricated prosthetic hand}
    (Chapter~4) and a \textbf{designed control application}
    (Chapter~11) --- an offline web app with a 3D digital twin,
    live mapping and pose editors, and a two-contract firmware
    interface --- built to the point of a working, hardware-verified
    firmware scaffold and a fully functional mock-mode web
    prototype.
    \item A \textbf{real-time inference pipeline} (Chapter~12) that
    reuses the offline-trained model unmodified, adds a
    deployment-appropriate fixed activity gate, confidence
    threshold, and \acrshort{fsm} confirmation layer, and meets its
    real-time latency budget with substantial margin (14\,ms per
    window against a 200\,ms step).
\end{itemize}

Consistent with the project scope defined in Chapter~1, this is
engineering proof-of-concept work: a single-subject,
discrete-command \acrshort{bci} system, not a clinically validated
or cross-subject-generalisable one. Within that scope, the system
meets its own design goals, and --- as important to this report as
the headline number itself --- every major design decision along
the way is backed by a documented experiment, including the ones
that failed (Chapters~6--9).

%% ----------------------------------------------------------
\section{Future Work}
\label{sec:conc_future}
%% ----------------------------------------------------------

Every item below was identified as concrete, unresolved future work
somewhere in this report; they are consolidated here by category
rather than repeated in narrative form.

\subsection{Data and Model}

\begin{enumerate}[leftmargin=*]
    \item \textbf{Clean single-blink re-recording.} The residual
    single-blink contamination (Chapters~6, 10, 13) is the largest
    remaining accuracy ceiling in the system and is not fixable by
    further classifier changes. A dedicated re-recording session
    with deliberate, experimenter-verified pauses of at least
    1.5\,s between individual blinks would directly test whether
    the \acrshort{dtw} architecture's remaining error is
    data-limited (as the current evidence strongly suggests,
    Chapter~10) rather than method-limited. The event-locked
    detection harness built during Attempt~4 of the blink
    investigation (Chapter~9) is already in place to validate any
    such re-recording.
    \item \textbf{Transfer learning across subjects.} The current
    system is single-subject by design (Chapters~1, 6). A natural
    extension is a general, multi-subject base model, fine-tuned per
    new user with a short (5--10 minute) personal calibration
    session, which would reduce the per-user data-collection burden
    that currently requires three full recording sessions
    (Chapter~6) before a new user can be supported at all.
    \item \textbf{Synchronised event markers.} The failed
    \code{Random/} mixed-gesture validation attempt (Chapter~6)
    traced its failure to the marker-logging application and the
    EEG amplifier running on two separate, unsynchronised clocks.
    Recording event markers directly on the EEG acquisition device
    itself, on the same clock as the signal, would resolve this and
    enable a genuine free-running, mixed-gesture validation set ---
    something the project currently lacks entirely.
\end{enumerate}

\subsection{Hardware Deployment}

\begin{enumerate}[leftmargin=*]
    \item \textbf{Emotiv EPOC~X validation.} As discussed in
    Chapter~13, every result in this report was produced on the
    19-channel data-collection device, not the 14-channel Emotiv
    EPOC~X that was always the intended deployment target. The
    approximate channel correspondence identified (T3/T4
    $\rightarrow$ T7/T8; Fp1/Fp2 $\rightarrow$ AF3/AF4) has not been
    validated with real recordings. The recommended path is direct
    re-calibration on the target device --- collecting a fresh
    multi-session dataset on the Emotiv itself, following the same
    protocol as Chapter~6 --- rather than attempting to transfer the
    existing model via the approximate channel mapping, since a
    transferred model's accuracy on different hardware, electrode
    geometry, and the Emotiv's onboard 45\,Hz hardware filter cannot
    be assumed without direct testing. Raw signal access on the
    Emotiv device (bypassing the vendor's Cortex software layer, via
    a library such as CyKit) would be needed to reproduce this
    project's preprocessing pipeline (Chapter~7) exactly.
    \item \textbf{Real-time activity gate recalibration.} The fixed
    20\,\si{\micro\volt} real-time activity gate (Chapter~12) was
    calibrated against amplitude distributions from the
    data-collection hardware and has not been validated on a
    different device or across multiple live sessions. A short,
    automatic per-session calibration routine --- for example,
    sampling a brief rest period at the start of each session to set
    a session-specific threshold, approximating the adaptivity of
    the offline recording-relative gate (Chapter~7) without
    requiring a full recording in advance --- would make the
    real-time gate more robust to exactly the kind of session-to-session
    variation \acrshort{loso} evaluation (Chapter~9) is designed to
    stress-test.
\end{enumerate}

\subsection{System Integration}

\begin{enumerate}[leftmargin=*]
    \item \textbf{Complete Firmware Layers 1--5.} Only Layer~0 is
    verified on physical hardware at the time of writing
    (Chapter~11). Servo control and fail-safe (Layer~1), the WiFi
    access point and live status dashboard (Layer~2), the mapping
    editor (Layer~3), the pose creator (Layer~4), and the 3D digital
    twin (Layer~5, stretch goal) all remain to be built and flashed.
    \item \textbf{Resolve the Contract~A label mismatch.} The
    shipped \code{contracts.h} and default \code{mappings.json}
    currently define a placeholder seven-label set that does not
    match the classifier's real five-label output (Chapter~11). This
    is a small, well-scoped fix (Table~\ref{tab:label_reconciliation})
    but is a blocking prerequisite for Layer~1, since the
    \acrshort{uart} listener would otherwise be listening for labels
    the classifier will never send.
    \item \textbf{End-to-end integration testing.} The real-time
    inference pipeline (Chapter~12) and the firmware
    (Chapter~11) have each been validated independently but never
    run together against physical hardware: the inference pipeline
    against recorded data and the offline-trained model, and the
    firmware only up to Layer~0. Connecting the two --- live
    \acrshort{eeg} through inference through \acrshort{uart} through
    the ESP32 to the physical servos --- is the natural next
    integration milestone once Firmware Layer~1 exists.
    \item \textbf{Implement saved gesture sequences.} Identified in
    Chapter~11 as a planned-but-unimplemented feature from an
    earlier project planning document: chaining multiple saved poses
    into an ordered, replayable sequence, exposed through a new
    Contract~B endpoint and an extended pose data model. This is a
    web-app and firmware feature addition, not a classification or
    signal-processing change.
\end{enumerate}

%% ----------------------------------------------------------
\section{Closing Remarks}
\label{sec:conc_remarks}
%% ----------------------------------------------------------

This report has tried to document not just what the final system
achieves, but how it got there --- including the evaluation crisis
that nearly went unnoticed, the four blink-classification methods
that failed before a fifth succeeded, and the honest, unresolved
gaps that remain in firmware integration and target-hardware
validation. That choice was deliberate: a graduation project's value
lies as much in the documented reasoning behind its design decisions
as in its headline accuracy figure, and a 91.3\% result reached via
a rigorous, leakage-audited protocol and a physiology-matched
architecture is worth more than a higher number reached by a less
careful evaluation.

As stated in Chapter~1, this project is explicitly scoped as
engineering proof-of-concept, not clinical validation: it does not
include trials with amputee subjects, formal usability studies, or
regulatory review. What it does provide is a complete, working,
end-to-end demonstration --- from a voluntary blink or jaw
contraction, through a rigorously evaluated classifier, to a
physical prosthetic hand actuating the corresponding command --- and
a clearly enumerated set of concrete next steps
(Section~\ref{sec:conc_future}) for turning that demonstration into
a system ready for the target consumer hardware and, eventually, for
the kind of user-facing evaluation that engineering proof-of-concept
work is meant to earn the right to attempt.